%==============================================================================
\section{Stochastic Gradient Descent (SGD)}
\label{sec:sgd}
%==============================================================================

\begin{dinhnghia}[title={Stochastic Gradient Descent (SGD)}]
Kỹ thuật tối ưu phổ biến trong học máy: cập nhật tham số (trọng số, bias) lặp lại bằng \textbf{từng mẫu huấn luyện} (hoặc mini-batch) thay vì toàn bộ dataset.
Biến thể của gradient descent; hiệu quả và nhanh hơn với dataset lớn và thưa.
\end{dinhnghia}

\subsection{Gradient Descent}

\begin{dinhnghia}[title={Gradient Descent}]
Thuật toán tối ưu tối thiểu hóa \textbf{hàm cost} của mô hình: tính gradient cost theo tham số, cập nhật tham số ngược hướng gradient.
\end{dinhnghia}

\subsection{SGD}

\begin{dinhnghia}[title={SGD}]
Biến thể cập nhật tham số sau \textbf{mỗi ví dụ} (hoặc mini-batch nhỏ), không đợi duyệt hết dataset mỗi bước.
Hội tụ nhanh hơn và cần ít bộ nhớ lưu toàn bộ dữ liệu hơn batch gradient descent.
\end{dinhnghia}

\begin{phuongphap}[title={Thuật toán SGD}]
\begin{enumerate}
  \item Khởi tạo tham số ngẫu nhiên.
  \item Mỗi epoch: xáo trộn dữ liệu huấn luyện.
  \item Với mỗi mẫu: tính gradient cost theo tham số; cập nhật tham số ngược gradient.
  \item Lặp đến hội tụ.
\end{enumerate}
\end{phuongphap}

\begin{phuongphap}[title={Cập nhật trọng số}]
\[
w := w - \alpha \, \nabla_w J(w; x_i, y_i)
\]
$x_i$, $y_i$: mẫu và nhãn; $\alpha$: learning rate; $J$: hàm loss; $\nabla_w J$: gradient của $J$ theo $w$.
\end{phuongphap}

\begin{luuy}[title={Khác batch gradient descent}]
SGD dùng \textbf{một} mẫu (hoặc mini-batch) để tính gradient; batch GD dùng \textbf{toàn bộ} dataset mỗi bước.
\end{luuy}

\subsection{Ví dụ Python --- SGDClassifier trên Iris}

\begin{vidu}[title={Triển khai đầy đủ}]
Dự đoán loài hoa Iris từ hai feature sepal width và sepal length:

\begin{lstlisting}
import numpy as np
from sklearn import datasets
from sklearn.linear_model import SGDClassifier
from sklearn.model_selection import train_test_split
from sklearn.preprocessing import StandardScaler
from sklearn import metrics

iris = datasets.load_iris()
X, y = iris.data[:, :2], iris.target

X_train, X_test, y_train, y_test = train_test_split(
    X, y, test_size=0.20, random_state=1)

scaler = StandardScaler().fit(X_train)
X_train = scaler.transform(X_train)
X_test = scaler.transform(X_test)

clf = SGDClassifier(alpha=0.001, max_iter=200)
clf.fit(X_train, y_train)

y_train_pred = clf.predict(X_train)
print("The Accuracy of SGD classifier is:",
      metrics.accuracy_score(y_train, y_train_pred) * 100)
\end{lstlisting}

\textbf{Output mẫu:}
\begin{verbatim}
The Accuracy of SGD classifier is: 73.33333333333333
\end{verbatim}
\end{vidu}

\subsection{Ứng dụng}

\begin{ynghia}[title={Ứng dụng SGD}]
SGD là kỹ thuật tối ưu, không phải mô hình ML hoàn chỉnh; dùng rộng rãi khi dữ liệu \textbf{thưa} --- đặc biệt phân loại văn bản, NLP.
Mở rộng tốt cho bài toán hàng chục nghìn mẫu và hàng chục nghìn feature.
\end{ynghia}

\subsection{Ưu điểm và thách thức}

\begin{tinhchat}[title={Ưu điểm}]
\begin{itemize}
  \item \textbf{Hiệu quả bộ nhớ}: xử lý theo batch nhỏ.
  \item \textbf{Hội tụ nhanh} trên dataset lớn so với batch GD.
  \item \textbf{Thoát cực tiểu cục bộ}: tính ngẫu nhiên giúp tìm nghiệm tốt hơn đôi khi.
\end{itemize}
\end{tinhchat}

\begin{luuy}[title={Thách thức}]
\begin{itemize}
  \item \textbf{Gradient nhiễu}: có thể làm chậm hội tụ.
  \item \textbf{Chọn learning rate}: quan trọng cho tối ưu ổn định.
  \item \textbf{Kích thước mini-batch}: ảnh hưởng tốc độ hội tụ và ổn định.
\end{itemize}
\end{luuy}
